\documentclass[10pt,a4paper]{article}
\usepackage[UTF8]{ctex}
\usepackage{amsmath,amssymb,geometry,xcolor,listings,tcolorbox,multicol,enumitem}
\geometry{margin=1.2cm,top=1cm,bottom=0.8cm}
\setlength{\parskip}{0pt}\setlength{\parsep}{0pt}\setlength{\itemsep}{0pt}
\setlist{nosep,leftmargin=*}

\lstset{language=C,basicstyle=\ttfamily\scriptsize,keywordstyle=\color{blue}\bfseries,
  numbers=none,frame=single,breaklines=true,showspaces=false,tabsize=2,
  aboveskip=2pt,belowskip=2pt}

\newtcolorbox{keybox}[1][]{colback=yellow!3!white,colframe=yellow!70!black,
  title=\textbf{考点},fonttitle=\bfseries\small,size=small,top=2pt,bottom=2pt,#1}
\newtcolorbox{bugbox}[1][]{colback=red!3!white,colframe=red!60!black,
  title=\textbf{易错点},fonttitle=\bfseries\small,size=small,top=2pt,bottom=2pt,#1}

\title{\textbf{优先队列：堆 (Heap)}}
\author{}\date{}

\begin{document}
\maketitle\vspace{-8pt}

\section*{堆的两条性质}

\textbf{1. 结构性质：} 完全二叉树。\\
\textbf{2. 堆序性质：} 最小堆 — 父 $\le$ 所有孩子；最大堆 — 父 $\ge$ 所有孩子。

\section*{数组存储（索引从 1 开始）}

\begin{center}
\begin{tabular}{|c|c|c|}
\hline
\textbf{关系} & \textbf{公式} & \textbf{例子} \\
\hline
Parent(i) & $\lfloor i/2 \rfloor$ & parent(4)=2 \\
Left(i) & $2i$ & left(2)=4 \\
Right(i) & $2i+1$ & right(2)=5 \\
\hline
\end{tabular}
\end{center}

\textbf{data[0] 作 sentinel}（存极小值 MinData），自动终止上滤循环。

\begin{keybox}
\textbf{这三个公式只用于数组实现的堆！} 不适用于 BST、链表等其他结构。\\
\textbf{注意区分：} 原地堆排序用索引从0开始（left=2i+1, right=2i+2）。
\end{keybox}

\begin{multicols}{2}

% ====== 复杂度 ======
\section*{5种实现的复杂度对比}

\begin{tabular}{|l|c|c|}
\hline
\textbf{实现} & \textbf{Insert} & \textbf{DeleteMin} \\
\hline
无序数组 & $\Theta(1)$ & $\Theta(N)$ \\
无序链表 & $\Theta(1)$ & $\Theta(N)$ \\
有序数组 & $O(N)$ & $\Theta(1)$ \\
有序链表 & $O(N)$ & $\Theta(1)$ \\
BST(平衡) & $O(\log N)$ & $O(\log N)$ \\
\textbf{二叉堆} & $\mathbf{O(\log N)}$ & $\mathbf{O(\log N)}$ \\
\hline
\end{tabular}

\section*{初始化}

\begin{lstlisting}
PriorityQueue Init(int Max) {
    PriorityQueue H = malloc(sizeof(HeapStruct));
    H->Elements = malloc((Max+1)*sizeof(ElementType));
    H->Capacity = Max; H->Size = 0;
    H->Elements[0] = MinData;  // sentinel
    return H;
}
\end{lstlisting}

% ====== Insert ======
\section*{Insert — 上滤 (Percolate Up)}

\textbf{思想：} 新元素放在末尾（保持完全二叉树），然后和父节点比，小的往上冒。

\begin{lstlisting}
void insert(Heap* h, int x) {
    if(h->size >= MAX-1) return;
    int i = ++(h->size);                  // 先+后取！
    // 父节点比x大就下沉，空穴上移
    while(i>1 && h->data[i/2] > x) {
        h->data[i] = h->data[i/2];        // 父下沉
        i = i/2;                          // i上移
    }
    h->data[i] = x;                       // 放好
}
// T(N) = O(log N)
\end{lstlisting}

\textbf{手算示例：} 初始堆 [3,6,5,9,10]，Insert(2)。\\
2放data[6] → 跟data[3]=5比，5>2，5下沉到6，i=3。\\
跟data[1]=3比，3>2，3下沉到3，i=1。到根，2放data[1]。\\
结果：[2,6,3,9,10,5]。

\begin{bugbox}
\texttt{int i = h->size++} 后自增拿的是旧值。必须 \texttt{++(h->size)}。\\
循环里写成 \texttt{data[i/2]=x} 做了交换 → 应该只下沉，最后才放。\\
忘了循环外的 \texttt{h->data[i]=x}。
\end{bugbox}

% ====== DeleteMin ======
\section*{DeleteMin — 下滤 (Percolate Down)}

\textbf{思想：} 取走根（最小），最后元素补到根，然后跟较小孩子比，大的往下沉。

\begin{lstlisting}
int deleteMin(Heap* h) {
    if(isEmpty(h)) return -1;
    int min = h->data[1];                 // 最小值
    int x = h->data[h->size--];           // 最后一个元素
    int i=1, child;
    while(2*i <= h->size) {               // 有左孩子
        child = 2*i;
        // 如果右孩子存在且更小，选右
        if(child < h->size && h->data[child+1] < h->data[child])
            child++;
        if(x <= h->data[child]) break;    // 放好了
        h->data[i] = h->data[child];      // 孩子上移
        i = child;
    }
    h->data[i] = x;                       // 放好
    return min;
}
// T(N) = O(log N)
\end{lstlisting}

\textbf{手算示例：} 堆 [2,6,3,9,10,5]，DeleteMin。\\
取走2，x=5。i=1，child=2(左) vs 3(右)，3更小。5>3，3上移到1，i=3。\\
i=3，child=6(左)，2*i=6 > size=5，无孩子，停。5放在3。\\
结果：[3,6,5,9,10]。

\begin{bugbox}
大小比较反了：最小堆找较小孩子，最大堆找较大孩子。\\
\texttt{child!=h->size} 条件漏了 → 可能访问越界。\\
忘了最后 \texttt{h->data[i]=x}。
\end{bugbox}

\columnbreak

% ====== 其他操作 ======
\section*{其他堆操作}

\begin{itemize}
\item \textbf{DecreaseKey(P, $\Delta$, H)：} 减小某键值 $\to$ 上滤。提高进程优先级。
\item \textbf{IncreaseKey(P, $\Delta$, H)：} 增大某键值 $\to$ 下滤。降低CPU占用。
\item \textbf{Delete(P, H)：} DecreaseKey(P, $\infty$, H) 使其到根 + DeleteMin(H)。
\item \textbf{BuildHeap(H)：} Floyd自底向上建堆 $O(N)$。
\end{itemize}

% ====== BuildHeap ======
\section*{BuildHeap — Floyd 法 $O(N)$}

\textbf{思想：} 从最后一个非叶节点(N/2)开始，往前逐个下滤。叶子不用处理。

\begin{lstlisting}
void percDown(Heap* h, int i) {
    int x = h->data[i], child;
    while(2*i <= h->size) {
        child = 2*i;
        if(child < h->size && h->data[child+1] < h->data[child])
            child++;
        if(x <= h->data[child]) break;
        h->data[i] = h->data[child];
        i = child;
    }
    h->data[i] = x;
}

void buildHeap(Heap* h, int arr[], int N) {
    for(int i=0; i<N; i++) h->data[i+1] = arr[i];
    h->size = N;
    for(int j = N/2; j >= 1; j--)          // 从最后一个非叶到根
        percDown(h, j);
}
// T(N) = O(N) — 不是 O(N log N)！
\end{lstlisting}

\begin{keybox}
\textbf{为什么 BuildHeap 是 $O(N)$ 不是 $O(N\log N)$？}

\textbf{定理：} 完美二叉树中所有节点高度和 = $2^{h+1}-1-(h+1) \approx N$。\\
大部分节点在底层(叶子)几乎不下沉，总工作量接近 $N$ 而非 $N\log N$。
\end{keybox}

\textbf{手算示例：} arr=[4,2,7,1,3,6,5], N=7。N/2=3。\\
下滤 data[3]=7：7 和 5 换 $\to$ [4,2,5,1,3,6,7]。\\
下滤 data[2]=2：跟孩子1和3比，1 \char`\< 2 $\to$ 2和1换 $\to$ [4,1,5,2,3,6,7]。\\
下滤 data[1]=4：跟1和5比，1 \char`\< 4 $\to$ 4和1换 $\to$ [1,4,5,2,3,6,7]。再跟2和3比，2 \char`\< 4 $\to$ 4和2换 $\to$ [1,2,5,4,3,6,7]。

\begin{bugbox}
\texttt{for(j=N/2; j>1; j--)} 漏了根 j=1。应为 \texttt{j>=1}。\\
拷数组写了 \texttt{arr+i}（指针地址）→ 应 \texttt{arr[i]}。\\
\texttt{h->size} 用 for 循环 ++ → 直接 \texttt{h->size=N}。\\
Floyd 写成 while(j!=1)+j/=2+percDown两下 → 应是简单 for。
\end{bugbox}

% ====== d-堆 ======
\section*{d-堆}

每个节点有 d 个孩子。DeleteMin 需 d-1 次比较选最小 → $O(d\log_d N)$。

\textbf{索引关系：} parent=$\lfloor(i+d-2)/d\rfloor$，\\
children= $d(i-1)+2, d(i-1)+3, \dots, di+1$。

\textbf{大d的优缺点：}
\begin{itemize}
\item [+] 树更矮，外存场景（磁盘I/O少）
\item [--] \texttt{*d / /d} 不是位运算（对比 \texttt{*2 / /2} 是移位）
\item [--] DeleteMin 比较次数 $d-1$ 增多
\end{itemize}

\section*{应用：找第k大元素}

\begin{tabular}{|l|c|}
\hline
\textbf{方法} & \textbf{复杂度} \\
\hline
排序取第k & $O(N\log N)$ \\
大小k的最小堆 & $O(N\log k)$ \\
BuildHeap + k次DeleteMax & $O(N + k\log N)$ \\
QuickSelect & 平均$O(N)$，最坏$O(N^2)$ \\
\hline
\end{tabular}

\end{multicols}
\end{document}
